\section{第一阶段可运行模型：把 CPU、MMIO、中断和 DMA 变成状态机}

前面几节把硬件事实映射到内核对象，本节再补最后一块：用一个可编译的 \cpp20 小模型，把“用户态不能访问内核页”“用户态不能访问 MMIO”“设备通过中断通知 CPU”“DMA 映射有生命周期”这些说法变成能失败、能断言、能复现的状态机。

这个模型不是为了模拟 x86-64 的全部细节。第一阶段需要的是最小硬件契约：只要这几个契约讲不清，后面的系统调用、调度、文件系统和网络 I/O 都会变成背概念。

\begin{keyidea}
一个适合原理卷的硬件模型不追求外设数量，而追求边界清楚：CPU 以某个模式发起访问，总线按地址把访问分发到 RAM 或 MMIO，内存权限决定访问是否成立，设备完成后投递中断，DMA engine 只能在映射有效期内写入被授权的 buffer。
\end{keyidea}

\subsection{模型覆盖的硬件契约}

\begin{longtable}{p{0.2\linewidth}p{0.36\linewidth}p{0.32\linewidth}}
\toprule
模块 & 模型里的结构 & 对应 OS 原理 \\
\midrule
CPU 特权级 & \code{PrivilegeLevel::Ring3} 与 \code{PrivilegeLevel::Ring0} & x86-64 CPL、系统调用入口、中断返回 \\
页权限 & \code{PagePermission} 与 \code{PhysicalMemory::check\_access} & 用户页、内核页、PTE 权限、保护异常 \\
总线解码 & \code{Bus::load8/store8} 按地址选择 RAM 或 MMIO & 物理地址空间、设备寄存器窗口、\code{ioremap} \\
MMIO 设备 & \code{SimpleDevice} 的 status/data/doorbell & 驱动寄存器协议、doorbell、副作用写 \\
中断控制器 & \code{InterruptController::raise/take\_next} & IRQ descriptor、pending 事件、handler 分发 \\
trap frame & \code{Cpu::handle\_one\_interrupt} 保存 RIP/RFLAGS/CPL & 内核入口、异常帧、返回用户态所需状态 \\
DMA 映射 & \code{DmaEngine::map/unmap/device\_write} & DMA API、IOMMU 授权、buffer 所有权 \\
\bottomrule
\end{longtable}

这张表里的每一行都刻意对应一个 OS 对象。CPU 模式对应“谁有资格执行特权操作”；页权限对应“用户指针为什么要检查”；MMIO 对应“设备为什么不能被当作普通变量”；中断对应“设备完成为什么能唤醒等待者”；DMA 对应“completion 前为什么不能释放 buffer”。

\subsection{状态机比名词更重要}

第一阶段最容易写成硬件名词表。名词表的问题是没有失败边界：读者知道 DMA、IOMMU、MMIO，却不知道什么时候该拒绝访问，什么时候该保存现场，什么时候资源已经不能再被设备使用。这个模型用几个状态转移强制表达边界。

\begin{enumerate}
\tightlist
\item 用户态写普通用户页：总线走 RAM 路径，页权限允许，写入成功。
\item 用户态写内核页：地址仍是 RAM，但页权限拒绝，抛出保护 fault。
\item 用户态写 MMIO doorbell：总线发现地址属于设备窗口，但模式不是 kernel，抛出保护 fault。
\item 内核写 MMIO doorbell：设备接收命令，更新内部状态，并向中断控制器投递 device interrupt。
\item CPU 处理 pending interrupt：保存 trap frame，切换到 Ring 0，把 RIP 改成设备中断入口。
\item 内核为 buffer 建立 DMA mapping：设备拿到 handle 后才能写入对应 RAM。
\item DMA unmap 后设备继续写：mapping 已失效，模型必须拒绝这次迟到写。
\end{enumerate}

这些状态转移对应真实内核里的不变量：权限检查必须发生在访问之前；中断入口必须保存足够恢复用户态的状态；设备拥有 buffer 的时间必须被内核精确管理；unmap 不是注释，而是权限撤销。

\subsection{编译和运行}

在书根目录运行下面两条命令，可以独立编译并执行本节模型。

\begin{lstlisting}
g++ -std=c++20 -Wall -Wextra -Wpedantic -Werror \
  source/latex/listings/stage_one_hardware_model.cpp \
  -o /tmp/stage_one_hardware_model
/tmp/stage_one_hardware_model
\end{lstlisting}

期望输出如下：

\begin{lstlisting}
stage one hardware model checks passed
\end{lstlisting}

如果读者把模型里的任一条边界删掉，例如允许用户态写 MMIO、允许用户态写内核页、或者允许 unmap 后继续 DMA 写，程序都应该失败。这个失败比“书上说不可以”更有教学价值，因为它把抽象规则变成了可观察的断言。

\subsection{测试断言和失败含义}

\begin{longtable}{p{0.3\linewidth}p{0.34\linewidth}p{0.24\linewidth}}
\toprule
断言 & 失败说明 & 对应真实系统风险 \\
\midrule
用户页读写成功 & 基础 RAM 路径或页权限初始化错误 & 用户程序无法正常执行 \\
用户态写内核页失败 & 缺少特权隔离 & 用户进程可破坏内核对象 \\
用户态写 MMIO 失败 & 设备寄存器暴露给用户态 & 用户程序可直接控制硬件 \\
内核写 doorbell 后产生中断 & 设备完成路径没有闭环 & 等待者可能永久睡眠 \\
中断保存 trap frame & 返回路径缺少现场 & 用户态恢复后状态错乱 \\
DMA 写入映射 buffer 成功 & buffer 授权或地址翻译错误 & I/O 数据无法进入内存 \\
unmap 后 DMA 写失败 & 生命周期撤销无效 & 迟到 DMA 可能写坏新对象 \\
\bottomrule
\end{longtable}

这组断言不是完整 OS 测试，但它覆盖第一阶段必须过关的最小边界。后面讲系统调用时，可以把“用户指针检查”接回页权限；讲调度时，可以把“中断入口保存现场”接回 trap frame；讲文件系统和网络时，可以把“请求完成”接回中断与 DMA 生命周期。

\subsection{完整 \cpp20 模型源码}

\lstinputlisting[language=C++,basicstyle=\ttfamily\tiny]{listings/stage_one_hardware_model.cpp}
